\section{\pimattn{} Design}

\subsection{System Overview}

\pimattn{} adopts a three-phase heterogeneous execution model.
Phase 1 (PIM): decomposed Q@K dot product on near-bank PIM with 2-bit KV.
Phase 2 (GPU): partial result reduction and softmax.
Phase 3 (GPU): score@V matrix multiply with dequantized float16 values.

\subsection{Decomposed Q@K on Near-Bank PIM}

For a query $Q \in \mathbb{R}^{M \times d}$ and key cache $K \in \mathbb{R}^{N \times d}$ striped across $N_{\text{banks}}$, each bank computes its partition of Eq.~\ref{eq:full}:

\begin{enumerate}
\item $\sum \hat{q}_j \hat{k}_{ij}$ --- 2-bit MAC, the main dot product.
\item $\sum \hat{q}_j$ --- 2-bit accumulate, computed once per query.
\item $\sum \hat{k}_{ij}$ --- 2-bit accumulate, per key vector.
\item Aggregate via Eq.~\ref{eq:full} after cross-bank reduction.
\end{enumerate}

Operations (1)--(3) all use 2-bit data, making them efficient on simple PIM PEs.
The final float16 scaling $s_Q s_K$ is applied after reduction.

\subsection{PIM Latency Model}

We model per-bank latency using a pipelined row-buffer execution model.
Data is loaded from DRAM into the row buffer ($R$ bytes), then the PE computes on it.
While PE computes on buffer $i$, buffer $i+1$ fills from DRAM.

Total data per GEMV: $B_{\text{total}} = (MK + KN + MN) \cdot e$ bytes, where $e$ is element size.
Per bank: $B_{\text{bank}} = B_{\text{total}} / N_{\text{banks}}$.
Row buffer fills: $n_{\text{rb}} = \lceil B_{\text{bank}} / R \rceil$, each taking $t_{\text{fill}} = R / B$ ns.
PE time: $T_{\text{pe}} = B_{\text{bank}} / W \cdot \tau$, where $\tau$ is the PE cycle time.

When asymmetric quantization is enabled, reduction terms add PE-only overhead:
\begin{equation}
T_{\text{red}} = \frac{(M+N) \cdot K \cdot e}{N_{\text{banks}} \cdot W} \cdot \tau
\end{equation}

The pipelined latency is:
\begin{equation}
T_{\text{latency}} = t_{\text{fill}} + \max\!\left((n_{\text{rb}}-1) \cdot t_{\text{fill}},\; T_{\text{pe}} + T_{\text{red}}\right)
\label{eq:pipeline}
\end{equation}

\subsection{Score@V on GPU}

After Q@K produces float16 scores $S$, the output is $O = S \cdot V$.
Since $S$ is float16 and $V$ must be dequantized from 2-bit, this operation is float16-intensive.
GPU provides 989 TFLOPS (B100) --- far exceeding per-bank PE throughput.
V dequantization ($v = s_V(\hat{v} - z_V)$) and matrix multiply are efficiently handled by GPU tensor cores.

\subsection{Heterogeneous Timeline}

GPU and PIM advance independent timelines:
\begin{equation}
T_{\text{total}} = T_{\text{PIM}}^{\text{QK}} + T_{\text{GPU}}^{\text{softmax}} + T_{\text{GPU}}^{\text{SV}}
\end{equation}
While Q@K, softmax, and score@V are sequential due to data dependencies, GPU projection operations (QKV projection, output projection) overlap with PIM attention execution.

\subsection{KV Cache Quantization Overhead}

In decode mode, each new token's K and V must be quantized before cache storage.
We model this on GPU:
\begin{equation}
T_{\text{quant}} = \max\!\left(\frac{2 \cdot d_{\text{kv}} \cdot e}{\text{BW}_{\text{GPU}}},\; \frac{4 \cdot d_{\text{kv}}}{\text{FLOPS}_{\text{GPU}}}\right)
\end{equation}
where $d_{\text{kv}} = n_{\text{kv\_heads}} \cdot d_{\text{head}}$.
For Llama-3-8B ($d_{\text{kv}}\!=\!1024$), this overhead is $<\!0.01\mu s$ per token --- negligible.
